\chapter{Laplasjan grafu, dyskretna teoria Hodge'a i twierdzenie Kirchhoffa}
\label{ch:tg-laplacian}

W rozdziale \ref{ch:tg-podstawy} pojawiła się macierz sąsiedztwa,
a w rozdziale \ref{ch:tg-algebra} macierz incydencji. Z ich połączenia otrzymamy operator, który
jest dyskretnym odpowiednikiem laplasjanu na rozmaitości.
W niniejszym rozdziale wykazujemy związek tego operatora
ze spójnością, potencjałem harmonicznym, rozkładem Hodge'a
oraz liczbą drzew rozpinających.

\section{Definicja i energia Dirichleta}
\label{sec:tg-lap-definition}

Niech $G=(V,E)$ będzie skończonym grafem prostym
o $n=|V|\geq1$ wierzchołkach, $m=|E|$ krawędziach
i $c$ składowych spójnych. Każdej krawędzi nadajmy dodatnią wagę $w_e>0$.
Macierz diagonalną tych wag oznaczamy przez $W$.

\begin{definicja}[Laplasjan ważony]
Niech $G=(V,E)$ będzie skończonym grafem prostym,
$n=|V|\geq1$, $m=|E|$, a $w_e>0$ dla $e\in E$.
Niech $W=\operatorname{diag}(w_e:e\in E)$ i
$B\in\mathbb R^{n\times m}$ będzie macierzą incydencji
grafu $G$ dla dowolnej orientacji pomocniczej. Definiujemy
\emph{laplasjan ważony} $L:=BWB^T:\mathbb R^V\to\mathbb R^V$.
Gdy wszystkie $w_e=1$, piszemy $L_G:=BB^T$ i mówimy
po prostu o laplasjanie grafu.
\end{definicja}

Dla grafu prostego nieważonego macierz ma postać
$L_G=D-A_G$, gdzie $D$ jest diagonalna i $D_{ii}=d_G(v_i)$.
Ważone elementy diagonalne są sumami wag krawędzi
przylegających do wierzchołka, a pozadiagonalny element
$L_{uv}$ wynosi $-w_{uv}$ dla krawędzi $uv$ i zero w pozostałych
przypadkach. Odwrócenie pomocniczej orientacji zmienia znak
kolumny $B$, ale nie zmienia iloczynu $BWB^T$.

\begin{uwaga}[Laplasjan kombinatoryczny i znormalizowany]
\label{rem:tg-laplacian-conventions}
Jeśli graf prosty $G$ jest $d$-regularny dla $d\geq1$,
to jego nieważony Laplasjan kombinatoryczny ma postać
$L_G=D-A_G=dI-A_G$. Natomiast Laplasjan znormalizowany
$\mathcal L=I-A_G/d$ z rozdziału \ref{ch:tg-expanders}
spełnia $\mathcal L=d^{-1}L_G$. Jeśli wartości własne obu macierzy numerujemy rosnąco
z krotnościami, to $\lambda_j(L_G)=d\lambda_j(\mathcal L)$
dla $j=1,\ldots,n$. W szczególności dla spójnego grafu o co najmniej
dwóch wierzchołkach spójność algebraiczna
z tego rozdziału jest $d$ razy większa od luki
znormalizowanej z rozdziału \ref{ch:tg-expanders}.
Dla grafu bez krawędzi ($d=0$) ten wzór normalizujący
nie jest określony; Laplasjan kombinatoryczny
$L_G=0$ pozostaje natomiast dobrze zdefiniowany.
\end{uwaga}

\begin{samepage}
\begin{twierdzenie}[Tożsamość energii]
\label{thm:tg-energy}
Niech $G=(V,E)$ będzie skończonym grafem prostym
o $n\geq1$ wierzchołkach i $c$ składowych spójnych,
z dodatnimi wagami krawędzi i laplasjanem ważonym $L$.
Dla dowolnego $f\in\mathbb R^V$ zachodzi
\[
f^TLf=\sum_{\{u,v\}\in E}w_{uv}(f(u)-f(v))^2\ge0.
\]
Stąd $L$ jest macierzą symetryczną dodatnio półokreśloną, a
$\ker L$ składa się dokładnie z funkcji stałych na każdej
składowej spójnej. W szczególności $\dim\ker L=c$ i
$\operatorname{rank}L=n-c$.
\end{twierdzenie}
\end{samepage}

\begin{proof}
Ponieważ $L=BWB^T$, dostajemy
$f^TLf=(B^Tf)^TW(B^Tf)$. Dla krawędzi skierowanej
pomocniczo $u\to v$ jej współrzędna wektora $B^Tf$
wynosi $f(v)-f(u)$. Po rozwinięciu iloczynu otrzymujemy
wskazaną sumę kwadratów. Jest ona nieujemna, gdyż wagi są
dodatnie. Jeżeli $Lf=0$, suma kwadratów znika.
Odwrotnie, jeśli suma znika, to z dodatniości wszystkich
wag wynika $B^Tf=0$, a stąd $Lf=BWB^Tf=0$.
Zatem $Lf=0$ dokładnie wtedy, gdy na końcach każdej
krawędzi $f$ przyjmuje równe wartości.
To oznacza dokładnie stałość na składowych spójnych.
Funkcje charakterystyczne $\chi_C$ składowych spójnych $C$ stanowią bazę jądra, więc jego wymiar wynosi
$c$, a rząd wynosi $n-c$.
\end{proof}

\begin{przyklad}[Laplasjan ścieżki]
Dla ścieżki $1-2-3$ bez dodatkowych wag otrzymujemy
\[
L=\begin{pmatrix}1&-1&0\\-1&2&-1\\0&-1&1\end{pmatrix},
\quad
f^TLf=(f_1-f_2)^2+(f_2-f_3)^2.
\]
Wektor stały $(1,1,1)^T$ należy do jądra.
Wektory $(1,0,-1)^T$ i $(1,-2,1)^T$ są wektorami własnymi
odpowiadającymi wartościom odpowiednio $1$ i $3$.
Wartości własne laplasjanu wynoszą więc $0,1,3$.
\end{przyklad}

\section{Harmoniczność, zasada maksimum i zagadnienie Dirichleta}
\label{sec:tg-harmonic}

\begin{definicja}[Funkcja harmoniczna we wierzchołku]
Niech $G=(V,E)$ będzie skończonym grafem prostym
z dodatnimi wagami i laplasjanem ważonym $L$.
Dla $f:V\to\mathbb R$ oraz $v\in V$ mówimy, że $f$ jest \emph{harmoniczna w wierzchołku $v$},
jeżeli $(Lf)(v)=0$. Oznaczmy przez $d_w(v):=\sum_{u:\,uv\in E}w_{uv}$
jego ważony stopień.
\end{definicja}

Jeżeli $v$ jest izolowany, to $(Lf)(v)=0$ dla każdej funkcji,
ale wzór średniej ważonej nie jest wtedy zdefiniowany.
Dalej stosujemy go wyłącznie dla $d_w(v)>0$.

Warunek harmoniczności w nieizolowanym $v$ jest równoważny równości
\[
f(v)=\frac{1}{d_w(v)}\sum_{u:\,uv\in E}w_{uv}f(u),
\]
to znaczy wartość w wierzchołku jest średnią ważoną wartości
na jego sąsiadach. Jest to dyskretny odpowiednik
własności wartości średniej funkcji harmonicznych.

\begin{twierdzenie}[Zasada maksimum na grafie]
Niech $G$ będzie skończonym spójnym grafem ważonym, którego
wszystkie wagi są dodatnie. Niech $\varnothing\ne S\subseteq V$
będzie zbiorem wierzchołków brzegowych. Jeżeli funkcja
$f:V\to\mathbb R$ jest harmoniczna w każdym wierzchołku
z $V\setminus S$, to osiąga maksimum i minimum na $S$.
\end{twierdzenie}

\begin{proof}
Niech $M=\max_{v\in V}f(v)$ i niech $U$ będzie zbiorem
wierzchołków, w których funkcja przyjmuje wartość $M$.
Jeżeli $U\cap S\ne\varnothing$, maksimum jest osiągnięte
na brzegu. W przeciwnym razie każdy $v\in U$ jest
wewnętrzny i harmoniczny. Równość średniej ważonej
$f(v)=\sum_{u\sim v}w_{uv}f(u)/d_w(v)=M$ przy dodatnich wagach
implikuje, że każdy sąsiad $u$ wierzchołka $v$ też ma
wartość $M$. Zatem $U$ jest zamknięty na sąsiadów.
Spójność grafu prowadzi wtedy do $U=V$, co jest
sprzeczne z $U\cap S=\varnothing$.
Tak więc maksimum jest osiągnięte na $S$.
Stosując ten sam argument do $-f$, otrzymujemy
analogiczne stwierdzenie dla minimum.
\end{proof}

\begin{twierdzenie}[Dyskretne zagadnienie Dirichleta]
\label{thm:tg-dirichlet}
Niech $G=(V,E)$ będzie skończonym grafem prostym
z dodatnimi wagami i laplasjanem ważonym $L$,
a $S\subseteq V$ niepustym zbiorem przecinającym
każdą składową spójną. Dla dowolnej funkcji brzegowej
$g:S\to\mathbb R$ istnieje dokładnie jedna funkcja
$f:V\to\mathbb R$ taka, że $f|_S=g$ oraz
$(Lf)(v)=0$ dla każdego $v\in V\setminus S$.
Ponadto $f$ jest jednoznacznym minimizatorem energii
$E(h):=\tfrac12 h^TLh$ wśród wszystkich $h$ o $h|_S=g$.
\end{twierdzenie}

\begin{proof}
Niech $I=V\setminus S$. Jeśli $I=\varnothing$, funkcja $f=g$
jest jedynym rozwiązaniem i twierdzenie jest oczywiste.
W przeciwnym razie rozpiszmy macierz $L$ według podziału
$V=I\sqcup S$ na bloki $L_{II},L_{IS},L_{SI},L_{SS}$.
Równanie $Lf=0$ na $I$ przy zadanych wartościach
$f_S=g$ ma postać $L_{II}f_I=-L_{IS}g$.

Pokażemy, że $L_{II}$ jest dodatnio określona. Niech
$x\in\mathbb R^I$, $x\ne0$, a $\widetilde x\in\mathbb R^V$
będzie jej przedłużeniem przez zero na $S$.
Wówczas $x^TL_{II}x=\widetilde x^TL\widetilde x\ge0$.
Równość oznaczałaby, na mocy twierdzenia
\ref{thm:tg-energy}, że $\widetilde x$ jest stała
na każdej składowej. Ponieważ każda składowa spotyka $S$,
a tam $\widetilde x=0$, wówczas $\widetilde x=0$,
sprzecznie z $x\ne0$. A więc $L_{II}$ jest dodatnio
określona i odwracalna; równanie ma jednoznaczne rozwiązanie.

Każda funkcja $h$ o tych samych wartościach na $S$
ma postać $h=f+\eta$ z $\eta|_S=0$.
Rozwinięcie energii daje
$E(f+\eta)=E(f)+\eta^TLf+\tfrac12\eta^TL\eta$.
Wyraz mieszany jest zerowy, ponieważ $Lf$ znika na
$I$, a $\eta$ na $S$. Pozostały składnik jest dodatni
dla $\eta\ne0$ na mocy dodatniej określoności $L_{II}$.
Otrzymujemy jednoznaczność minimum.
\end{proof}

\begin{uwaga}[Konieczność warunku na zbiór brzegowy]
Jeśli pewna składowa spójna $C$ grafu $G$ nie spotyka
zbioru $S$, to $\chi_C$ należy do jądra $L$ i znika na $S$.
Dodanie do dowolnego rozwiązania wielokrotności
$\chi_C$ zachowuje dane brzegowe i harmoniczność,
a także energię. Rozwiązanie nadal istnieje: na
składowych spotykających $S$ stosujemy twierdzenie
\ref{thm:tg-dirichlet}, na pozostałych wybieramy funkcje
stałe. Wszelkie różnice rozwiązań są stałe na każdej
składowej i zerowe na składowych spotykających $S$.
Istotnie, dla różnicy $h$ mamy $h|_S=0$ oraz
$(Lh)|_{V\setminus S}=0$, więc $h^TLh=0$.
Z tożsamości energii wynika żądana charakterystyka.
W szczególności wymiar przestrzeni różnic jest równy
liczbie składowych nieprzecinających $S$.
\end{uwaga}

\begin{przyklad}[Napięcie na łańcuchu trzech wierzchołków]
Na ścieżce $a-b-c$ ustalamy $f(a)=0$, $f(c)=1$
i przyjmujemy wagi jednostkowe. Harmoniczność w $b$ daje
$2f(b)-f(a)-f(c)=0$, więc $f(b)=1/2$.
Energia wynosi $E(f)=\tfrac12((1/2)^2+(1/2)^2)=1/4$.
Każda inna wartość $f(b)=t$ daje
$E(t)=\tfrac12(t^2+(1-t)^2)=t^2-t+1/2$,
które przyjmuje minimum dokładnie dla $t=1/2$.
\end{przyklad}

\begin{center}
\begin{tikzpicture}[scale=1.05,line cap=round]
 \node[tgBlueVertex] (a) at (-2.1,0) {$0$};
 \node[tgTealVertex] (b) at (0,0) {$\frac12$};
 \node[tgRedVertex] (c) at (2.1,0) {$1$};
 \draw[tgTree,line width=3pt] (a)--(b);
 \draw[tgAccent,line width=3pt] (b)--(c);
 \node[text=tgNavy,font=\footnotesize] at (0,-.67)
   {Wartość harmoniczna w środku to średnia wartości na końcach.};
\end{tikzpicture}
\end{center}

\section{Dyskretne operatory Hodge'a}
\label{sec:tg-hodge-operators}

Dla przestrzeni euklidesowych $C^0(G):=\mathbb R^V$
i $C^1(G):=\mathbb R^E$ wprowadźmy operator
różniczki dyskretnej $d_0:=B^T:C^0\to C^1$.
Dla krawędzi $e:u\to v$ otrzymujemy
$(d_0f)(e)=f(v)-f(u)$.
Operator sprzężony wynosi $d_0^*=B:C^1\to C^0$.

\begin{definicja}[Laplasjany Hodge'a na grafie]
Niech $G$ będzie skończonym grafem prostym z jednostkowymi
wagami i macierzą incydencji $B$. Definiujemy operator
na $0$-kołańcuchach
$\Delta_0:=d_0^*d_0=BB^T$ i operator na
$1$-kołańcuchach $\Delta_1:=d_0d_0^*=B^TB$.
Dla $k\in\{0,1\}$ kołańcuch $\omega\in C^k(G)$
nazywamy harmonicznym,
gdy $\Delta_k\omega=0$.
\end{definicja}

\begin{samepage}
\begin{twierdzenie}[Harmoniczne reprezentanty kohomologii]
\label{thm:tg-hodge-harmonic}
Niech $G$ będzie skończonym grafem prostym o $n$ wierzchołkach,
$m$ krawędziach i $c$ składowych. Wtedy
\[
\ker\Delta_0=\ker B^T\cong\mathbb R^c,
\qquad
\ker\Delta_1=\ker B\cong\mathbb R^{m-n+c}.
\]
Ponadto każda klasa w
$H^1(G;\mathbb R):=C^1(G)/\operatorname{im}d_0$
ma dokładnie jednego przedstawiciela należącego do
$\ker\Delta_1$. Ponieważ realizacja grafu jest kompleksem CW o wymiarze
co najwyżej jeden, operator $d_1:C^1\to C^2=0$ jest zerowy,
a więc $H^1=\ker d_1/\operatorname{im}d_0$.
\end{twierdzenie}
\end{samepage}

\begin{proof}
Ponieważ $\langle\Delta_0f,f\rangle=\|B^Tf\|^2$,
warunek $\Delta_0f=0$ jest równoważny $B^Tf=0$.
Oznacza to stałość $f$ na każdej składowej,
a więc wymiar jądra wynosi $c$.
Analogicznie $\langle\Delta_1h,h\rangle=\|Bh\|^2$,
zatem $\ker\Delta_1=\ker B$; jego wymiar
wynosi $m-n+c$ z twierdzenia \ref{thm:tg-rank}.

Z rozkładu ortogonalnego
$C^1=\operatorname{im}B^T\oplus^\perp\ker B$
każdy $x\in C^1$ rozkłada się jednoznacznie na
$x=d_0f+h$, gdzie $h\in\ker B$.
Jednoznaczne są składniki $d_0f$ i $h$, natomiast
$f$ jest wyznaczona z dokładnością do funkcji stałej
na każdej składowej spójnej. W klasie $[x]$
modulo obrazu $d_0$ przedstawiciel $h$ istnieje,
a gdyby były dwa, ich różnica należałaby zarówno
do $\operatorname{im}B^T$, jak i do jego dopełnienia
ortogonalnego. Musiałaby więc być zerowa.
\end{proof}

\begin{uwaga}[Porównanie z rozmaitościami]
W teorii Hodge'a na rozmaitości zewnętrzna różniczka
ma postać $d:\Omega^k\to\Omega^{k+1}$, a laplasjan
Hodge'a to $\Delta=dd^*+d^*d$. Na grafie mamy tylko
wymiary zero i jeden, więc $d_1=0$: stąd na
$C^0$ pozostaje $d_0^*d_0$, a na $C^1$ -- $d_0d_0^*$.
To skończeniewymiarowy odpowiednik twierdzenia Hodge'a
dla zwartych orientowalnych rozmaitości riemannowskich
bez brzegu, gdzie każda klasa kohomologii de Rhama
ma jednoznacznego harmonicznego reprezentanta.
\end{uwaga}

\begin{stwierdzenie}[Dodatnie widma dwóch laplasjanów]
Niech $B$ będzie rzeczywistą macierzą incydencji
skończonego grafu prostego.
Dodatnie wartości własne macierzy $BB^T$ i $B^TB$
są takie same, z tymi samymi krotnościami.
\end{stwierdzenie}

\begin{proof}
Jeśli $BB^Tf=\lambda f$ z $\lambda>0$, to
$B^Tf\ne0$ (gdyby $B^Tf=0$, lewa strona byłaby zerowa),
a ponadto
$B^TB(B^Tf)=B^T(BB^Tf)=\lambda B^Tf$.
W drugą stronę, dla $B^TBh=\lambda h$ wektor
$Bh$ jest niezerowy i spełnia
$BB^T(Bh)=\lambda Bh$.
Ograniczenia $B^T$ oraz $\lambda^{-1}B$ do odpowiednich
przestrzeni własnych są wzajemnie odwrotne, ponieważ
$BB^Tf=\lambda f$ oraz $B^TBh=\lambda h$.
Obie macierze są symetryczne, więc krotności
algebraiczne ich wartości własnych są równe wymiarom
odpowiednich przestrzeni własnych.
\end{proof}

\section{Twierdzenie macierzowo-drzewowe Kirchhoffa}
\label{sec:tg-matrix-tree}

\begin{definicja}[Liczba drzew rozpinających]
Dla skończonego grafu prostego $G$ o niepustym
zbiorze wierzchołków przez $\tau(G)$ oznaczamy liczbę jego drzew rozpinających.
Jeśli $G$ jest niespójny, przyjmujemy $\tau(G)=0$.
Dla grafu jednoelementowego $K_1$ przyjmujemy
$\tau(K_1)=1$, ponieważ jego jedynym drzewem
rozpinającym jest on sam.
\end{definicja}

\begin{uwaga}[Dwa różne niezmienniki]
W rozdziale \ref{ch:tg-skojarzenia} symbol
$\tau_V(G)$ oznacza najmniejszą moc pokrycia
wierzchołkowego. W tym rozdziale i w dalszych
zastosowaniach twierdzenia Kirchhoffa
$\tau(G)$ oznacza natomiast liczbę drzew
rozpinających. Nie należy utożsamiać
tych dwóch wielkości.
\end{uwaga}

\begin{twierdzenie}[Kirchhoff: liczba drzew z minora laplasjanu]
\label{thm:tg-kirchhoff}
Niech $G=(V,E)$ będzie skończonym spójnym grafem prostym,
$|V|=n\ge1$, z dodatnimi wagami $w_e$ na krawędziach.
Niech $B$ będzie macierzą incydencji, $W$ diagonalną
macierzą wag, $L=BWB^T$ i $r\in V$. Niech $L^{(r)}$
będzie macierzą otrzymaną
przez skreślenie wiersza i kolumny odpowiadających
wierzchołkowi $r\in V$. Wówczas
\[
\det L^{(r)}=\sum_{T}\prod_{e\in E(T)}w_e,
\]
gdzie suma biegnie po wszystkich drzewach rozpinających $T$.
Dla wag jednostkowych otrzymujemy
$\det L^{(r)}=\tau(G)$, niezależnie od wyboru $r$.
Dla $n=1$ wyznacznik pustej macierzy przyjmujemy równy $1$.
\end{twierdzenie}

\begin{proof}
Przypadek $n=1$ jest natychmiastowy. Niech $n\ge2$.
Usunięcie $r$-tego wiersza z macierzy $B$ daje
$(n-1)\times m$ macierz $B^{(r)}$. Mamy
$L^{(r)}=B^{(r)}W(B^{(r)})^T$.
Z wzoru Cauchy'ego--Bineta na wyznacznik iloczynu
macierzy prostokątnych otrzymujemy
\[
\det L^{(r)}=\sum_{\substack{S\subseteq E\\|S|=n-1}}
  \det(B^{(r)}_S)^2\prod_{e\in S}w_e,
\]
gdzie $B^{(r)}_S$ oznacza kwadratową podmacierz
złożoną z kolumn oznaczonych przez $S$.

Udowodnimy, że $\det(B^{(r)}_S)=0$, jeżeli $(V,S)$
nie jest drzewem, a $\det(B^{(r)}_S)=\pm1$, gdy jest
drzewem. Jeśli $(V,S)$ zawiera cykl, kolumny
incydencji tego cyklu są liniowo zależne: po wybraniu
znaków zgodnych z jego przebiegiem ich suma wynosi zero.
Zatem wyznacznik znika. Jeśli zaś $(V,S)$ nie zawiera
cyklu i ma $n-1$ krawędzi, jest lasem z
$n-(n-1)=1$ składową, czyli drzewem.

Niech więc $(V,S)$ będzie drzewem. W drzewie z co najmniej
dwoma wierzchołkami istnieje liść $v\ne r$.
W macierzy $B^{(r)}_S$ wiersz odpowiadający $v$
ma dokładnie jeden niezerowy element, równy $1$ albo $-1$.
Rozwinięcie wyznacznika względem tego wiersza sprowadza
go, z dokładnością do znaku, do analogicznego wyznacznika
dla drzewa z usuniętym liściem $v$ i przyległą krawędzią.
Powtarzając procedurę, dochodzimy do pustej macierzy
o wyznaczniku $1$. Stąd $\det(B^{(r)}_S)=\pm1$.

Do sumy Cauchy'ego--Bineta wkład dają więc dokładnie
zbiory krawędzi tworzące drzewa rozpinające, a ich
wkład wynosi $\prod_{e\in E(T)}w_e$.
Po zsumowaniu otrzymujemy tezę.
Warto zauważyć, że argument nie wymaga w istocie dodatniości
wag: jest to tożsamość wielomianowa w zmiennych $(w_e)_{e\in E}$.
Założenie $w_e>0$ zapewnia dodatkowo dodatniość każdego
składnika sumy i nieosobliwość zredukowanego laplasjanu
spójnego grafu. Sam wzór wyznacznikowy obowiązuje również
dla wag zerowych i ujemnych.
\end{proof}

\begin{przyklad}[Cztery wierzchołki na cyklu]
Dla cyklu $C_4$ z krawędziami jednostkowymi mamy
\[
L=\begin{pmatrix}
2&-1&0&-1\\-1&2&-1&0\\0&-1&2&-1\\-1&0&-1&2
\end{pmatrix}.
\]
Skreślając czwarty wiersz i kolumnę, dostajemy
\[
\det\begin{pmatrix}2&-1&0\\-1&2&-1\\0&-1&2\end{pmatrix}=4.
\]
Z drugiej strony każde drzewo rozpinające $C_4$
powstaje przez usunięcie dokładnie jednej z czterech
krawędzi. Takich drzew jest cztery, więc oba rachunki
się zgadzają.
\end{przyklad}

\begin{przyklad}[Ważony trójkąt]
Niech wagi trzech krawędzi trójkąta o wierzchołkach
$1,2,3$ wynoszą $w_{12}=2$, $w_{23}=w_{31}=1$.
Wtedy
\[
L=\begin{pmatrix}3&-2&-1\\-2&3&-1\\-1&-1&2\end{pmatrix},
\qquad
\det L^{(3)}=\det\begin{pmatrix}3&-2\\-2&3\end{pmatrix}=5.
\]
Trzy drzewa rozpinające składają się z dwóch spośród
trzech krawędzi, a ich iloczyny wag wynoszą $2$, $2$
i $1$. Suma $2+2+1=5$ zgadza się z wynikiem macierzowym.
\end{przyklad}

\begin{center}
\begin{tikzpicture}[scale=.95,line cap=round]
  \node[tgBlueVertex] (a) at (-1.7,-.95) {$1$};
  \node[tgBlueVertex] (b) at (1.7,-.95) {$2$};
  \node[tgBlueVertex] (c) at (0,1.40) {$3$};
  \draw[tgAccent,line width=3.0pt] (a)--node[below,font=\small] {$2$}(b);
  \draw[tgTree,line width=2.6pt] (b)--node[above right,font=\small] {$1$}(c);
  \draw[tgBlue,line width=2.6pt] (c)--node[above left,font=\small] {$1$}(a);
  \node[text=tgNavy,font=\footnotesize] at (0,-1.65)
    {Waga drzewa = iloczyn wag jego krawędzi.};
\end{tikzpicture}
\end{center}

\section{Minory główne laplasjanu i lasy ukorzenione}
\label{sec:tg-rooted-forests}

Usunięcie jednego wiersza i jednej kolumny w twierdzeniu
Kirchhoffa odpowiada wybraniu korzenia drzewa.
Usunięcie większej liczby wierszy i odpowiadających im
kolumn prowadzi do lasów, w których każda składowa
zawiera dokładnie jeden z zadanych korzeni.

\begin{twierdzenie}[Wzór na lasy z zadanym zbiorem korzeni]
\label{thm:tg-rooted-forest-minor}
Niech $G=(V,E)$ będzie skończonym grafem prostym
o $n\geq1$ wierzchołkach z dodatnimi wagami $w_e$,
a $L$ jego laplasjanem ważonym, zdefiniowanym przez
macierz incydencji $B$ i macierz wag $W$ jako $BWB^T$.
Dla dowolnego niepustego zbioru $S\subseteq V$ niech
$I=V\setminus S$, a $L_{II}$ oznacza podmacierz
o wierszach i kolumnach indeksowanych przez $I$. Wówczas
\[
\det L_{II}=\sum_{F\in\mathcal F_S}\prod_{e\in E(F)}w_e,
\]
gdzie $\mathcal F_S$ jest rodziną lasów rozpinających $G$,
w których każda składowa zawiera dokładnie jeden
wierzchołek zbioru $S$. Przyjmujemy, że wyznacznik
pustej macierzy i pusty iloczyn wynoszą $1$.
Suma po pustej rodzinie wynosi $0$.
\end{twierdzenie}

\begin{proof}
Oznaczmy $k=|I|$. Dla $k=0$ jedynym dopuszczalnym
lasem jest graf na $V$ bez krawędzi, więc obie strony
wynoszą $1$. Dla $k>0$ niech $B_I$ będzie macierzą
powstałą przez pozostawienie wierszy $B$ odpowiadających
$I$. Jeżeli $m=|E|<k$, rząd $L_{II}$ jest mniejszy
od $k$ i nie istnieje dopuszczalny las, który musiałby
mieć $k$ krawędzi. Obie strony są więc zerowe.
W przeciwnym przypadku wzór Cauchy'ego--Bineta daje
\[
\det L_{II}=\sum_{\substack{A\subseteq E\\|A|=k}}
    \det(B_{I,A})^2\prod_{e\in A}w_e.
\]
Jeżeli $(V,A)$ zawiera cykl, odpowiadające mu kolumny
macierzy $B$ są liniowo zależne, a ta zależność zachodzi
również po ograniczeniu do wierszy $I$.
Taki składnik ma zerowy wyznacznik.

Jeżeli $(V,A)$ jest lasem, to liczba jego składowych
wynosi $n-k=|S|$. Gdy pewna składowa $C$ nie zawiera
żadnego elementu $S$, wszystkie jej wierzchołki należą
do $I$. Suma odpowiadających im wierszy $B_{I,A}$
jest zerowa: krawędzie lasu albo mają oba końce
w $C$, albo żaden. Jest to nietrywialna zależność
wierszy. Jeżeli natomiast każda składowa spotyka $S$,
to z równości liczby składowych i liczby korzeni
wynika, że każda zawiera dokładnie jeden korzeń.

W tym ostatnim przypadku porządkujemy wiersze
i kolumny według składowych. Powstaje macierz
blokowo diagonalna, której blok dla składowej $C$
jest macierzą incydencji drzewa $C$ z usuniętym
wierszem jego korzenia. Każdy blok jest kwadratowy,
a jego wyznacznik wynosi $\pm1$ na mocy argumentu
z usuwaniem liści z dowodu twierdzenia
\ref{thm:tg-kirchhoff}. Składowa jednowierzchołkowa
odpowiada pustemu blokowi. Permutowanie wierszy
i kolumn może zmienić tylko znak wyznacznika,
zatem $\det(B_{I,A})^2=1$.
Po podstawieniu do sumy otrzymujemy tezę.
\end{proof}

\begin{uwaga}[Kiedy minor Dirichleta jest dodatni?]
Przy dodatnich wagach wzór
\ref{thm:tg-rooted-forest-minor} pokazuje, że
$\det L_{II}>0$ dokładnie wtedy, gdy $S$
spotyka każdą składową spójną grafu.
Istnienie dopuszczalnego lasu oczywiście wymaga
tego warunku. Aby wykazać wystarczalność, dla
każdego $v\notin S$ wybieramy sąsiada $p(v)$
o odległości od zbioru $S$ mniejszej o jeden,
korzystając z najkrótszej drogi do $S$.
Wybieramy wszystkie krawędzie $\{v,p(v)\}$.
Są parami różne: ich ponowny wybór przez drugi
koniec wymagałby jednoczesnego spadku odległości
w obu kierunkach. Nie tworzą cyklu, bo wierzchołek
o największej odległości od $S$ na hipotetycznym
cyklu musiałby wybrać obie krawędzie cyklu
(krawędzie łączą różne poziomy), a wybiera tylko jedną.
Każdy wierzchołek łączy się z pewnym korzeniem
przez kolejne przejścia do $p(v)$. Las ma
$n-|S|$ krawędzi, a więc $|S|$ składowych,
z których każda zawiera co najmniej jeden korzeń.
Zatem każda zawiera dokładnie jeden.
Dodatniość minora wynika z istnienia takiego lasu.
Jest to kombinatoryczna kontrola warunku
odwracalności macierzy z twierdzenia
\ref{thm:tg-dirichlet}, także dla $I=\varnothing$.
\end{uwaga}

\begin{wniosek}[Wielomian lasów ukorzenionych]
\label{cor:tg-rooted-forest-polynomial}
Niech $G=(V,E)$ będzie skończonym grafem prostym
o $n\geq1$ wierzchołkach, z dodatnimi wagami $w_e$
i laplasjanem ważonym $L$. Dla $t\in\mathbb R$
zachodzi
\[
\det(L+tI_n)=
\sum_F t^{c(F)}
\left(\prod_{C\in\operatorname{Comp}(F)}|V(C)|\right)
\prod_{e\in E(F)}w_e,
\]
gdzie suma biegnie po wszystkich lasach rozpinających
$F$ grafu $G$, $\operatorname{Comp}(F)$ jest zbiorem ich
składowych, a $c(F)$ oznacza liczbę tych składowych.
\end{wniosek}
\begin{proof}
Z wieloliniowości wyznacznika względem kolumn
otrzymujemy
\[
\det(L+tI_n)=\sum_{S\subseteq V}t^{|S|}
\det L_{V\setminus S,V\setminus S}.
\]
Istotnie, dla kolumn o indeksach w $S$ wybieramy $t$ razy
wektor bazy standardowej, a rozwinięcie względem
tych kolumn usuwa odpowiadające im wiersze i kolumny
bez zmiany znaku, ponieważ usuwamy te same indeksy
wierszy i kolumn.

Składnik $S=\varnothing$ znika,
ponieważ $L$ ma niezerowy wektor stały w jądrze.
Dla pozostałych składników stosujemy twierdzenie
\ref{thm:tg-rooted-forest-minor}.
Ustalony las $F$ występuje dla każdego zbioru $S$
zawierającego dokładnie jeden wierzchołek z każdej
składowej lasu. Takich zbiorów jest
$\prod_{C\in\operatorname{Comp}(F)}|V(C)|$ i wszystkie
mają moc $c(F)$. Grupowanie wyrazów daje wzór.
\end{proof}

\begin{przyklad}[Lasy ukorzenione na ścieżce]
Na ścieżce $a-b-c$ o wagach jednostkowych, przy
$S=\{a,c\}$ mamy $I=\{b\}$ i $\det L_{II}=2$.
Dopuszczalne są dokładnie dwa lasy: jeden zawiera
krawędź $ab$, drugi $bc$. Dla wszystkich zbiorów
korzeni otrzymujemy
$\det(L+tI_3)=t(t+1)(t+3)=t^3+4t^2+3t$.
Współczynnik przy $t$ pochodzi z jednego drzewa
i trzech wyborów korzenia. Współczynnik przy $t^2$
pochodzi z dwóch lasów, z których każdy ma
składowe o mocach $1$ i $2$. Współczynnik przy $t^3$
odpowiada lasowi bez krawędzi.
\end{przyklad}

\section{Wartości własne laplasjanu i liczba drzew}
\label{sec:tg-spectrum-tree}

\begin{twierdzenie}[Spektralny wzór Kirchhoffa]
\label{thm:tg-spectral-tree}
Niech $G$ będzie spójnym grafem prostym o $n\ge2$
wierzchołkach, z jednostkowymi wagami, a
$0=\lambda_1<\lambda_2\le\cdots\le\lambda_n$
niech będą wartościami własnymi jego laplasjanu
liczonymi z krotnościami. Wtedy
\[
\tau(G)=\frac1n\prod_{j=2}^n\lambda_j.
\]
\end{twierdzenie}

\begin{proof}
Z twierdzenia o energii i spójności jądro $L$ jest
jednowymiarowe, rozpięte przez $\mathbf1=(1,\ldots,1)^T$.
Z twierdzenia Kirchhoffa każdy główny kofaktor $L$
wynosi $\tau(G)$. Macierz dołączona
$\operatorname{adj}(L)$ (transpozycja macierzy dopełnień
algebraicznych) ma wszystkie kolumny w $\ker L$,
ponieważ $L\operatorname{adj}(L)=(\det L)I=0$.
Jej $j$-ta kolumna jest więc wielokrotnością
$\mathbf1$; element diagonalny tej kolumny jest
równy $\tau(G)$, co wyznacza tę wielokrotność.
Stąd $\operatorname{adj}(L)=\tau(G)\mathbf1\mathbf1^T$.

Rozważmy macierz
$M=L+\tfrac1n\mathbf1\mathbf1^T$.
Na prostej rozpiętej przez $\mathbf1$ działa ona jako
identyczność, a na dopełnieniu ortogonalnym jak $L$.
Zatem jej wartości własne, z uwzględnieniem krotności,
to $1,\lambda_2,\ldots,\lambda_n$
i $\det M=\prod_{j=2}^n\lambda_j$.
Z drugiej strony korzystamy z tożsamości
$\det(A+uv^T)=\det A+v^T\operatorname{adj}(A)u$,
ważnej również dla macierzy osobliwych.
Wynika ona z wieloliniowości wyznacznika względem
kolumn: składniki z co najmniej dwiema kolumnami
pochodzącymi z $uv^T$ znikają, a składniki z jedną
taką kolumną dają rozwinięcia według dopełnień
algebraicznych. Otrzymujemy
\[
\det M=\det L+\frac1n\mathbf1^T\operatorname{adj}(L)\mathbf1
=\frac1n\tau(G)n^2=n\tau(G).
\]
Porównanie obydwu wyrażeń kończy dowód.
\end{proof}

\begin{wniosek}[Wzór Cayleya]
Dla grafu pełnego $K_n$, $n\ge2$, liczba drzew
rozpinających wynosi $\tau(K_n)=n^{n-2}$.
\end{wniosek}

\begin{proof}
Macierz sąsiedztwa $K_n$ ma wszystkie elementy
pozadiagonalne równe $1$, natomiast stopnie wynoszą
$n-1$. Stąd $L=nI-J$, gdzie $J=\mathbf1\mathbf1^T$.
Na wektorze $\mathbf1$ operator znika, a na
$\mathbf1^\perp$ działa jako mnożenie przez $n$.
Zatem wartości własne to $0$ oraz $n$ o krotności $n-1$.
Ze wzoru spektralnego otrzymujemy
$\tau(K_n)=\frac1n n^{n-1}=n^{n-2}$.
\end{proof}

\begin{przyklad}[Widmo cyklu i grafu z przekątną]
Ponumerujmy wierzchołki $C_4$ kolejno od $1$ do $4$.
Wtedy
\[
L=\begin{pmatrix}
2&-1&0&-1\\-1&2&-1&0\\0&-1&2&-1\\-1&0&-1&2
\end{pmatrix}.
\]
Bezpośrednie mnożenie pokazuje, że wektory
\[
\begin{aligned}
u_0&=(1,1,1,1)^T,&u_1&=(1,0,-1,0)^T,\\
u_2&=(0,1,0,-1)^T,&u_3&=(1,-1,1,-1)^T
\end{aligned}
\]
są wektorami własnymi dla wartości odpowiednio
$0,2,2,4$. Są niezerowe i parami ortogonalne,
więc tworzą bazę $\mathbb R^4$.
Otrzymujemy $\tau(C_4)=\frac14(2\cdot2\cdot4)=4$.
Dodajmy przekątną $13$ o wadze jednostkowej.
Otrzymamy graf $K_4$ bez jednej krawędzi.
Jeśli $b=(1,0,-1,0)^T=u_1$, to nowy laplasjan
jest równy $L+bb^T$, na mocy wzoru $BWB^T$.
Ponieważ $b^Tu_1=2$ i $b^Tu_j=0$ dla $j\ne1$,
jedynie wartość własna odpowiadająca $u_1$
wzrasta z $2$ do $4$. Nowe widmo to $0,2,4,4$.
Stąd $\tau=\frac14(2\cdot4\cdot4)=8$.
Dodanie pojedynczej krawędzi podwoiło więc liczbę
rozpinających drzew w tym przykładzie.
\end{przyklad}

\section{Efektywna rezystancja i zasada minimalnej energii}
\label{sec:tg-resistance}

Laplasjan może również opisywać sieć elektryczną: wagę $w_e>0$
interpretujemy jako przewodność krawędzi, a jej odwrotność
$r_e=1/w_e$ jako opór. Ustalenie kierunków krawędzi jest wyłącznie
pomocnicze; to równania przepływu, a nie te kierunki, wyznaczają
prąd w sieci. W przyjętej konwencji macierzy $B$
kolumna krawędzi $u\to v$ ma znak $-1$ przy $u$
i $+1$ przy $v$. Niech $i_e$ oznacza prąd fizyczny, dodatni w kierunku
pomocniczym $u\to v$. Z prawa Ohma
$i_e=w_e(f(u)-f(v))$. W przyjętej konwencji
$Bi$ jest bilansem wpływającego prądu pomniejszonego
o wypływający. Dla jednostkowego źródła w $s$
i ujścia w $t$ mamy zatem $-Bi=[s]-[t]$, gdzie
$[v]$ oznacza wektor bazy standardowej przy $v$.
W dalszym ciągu zapisujemy $j=-i=WB^Tf$,
więc równania przyjmują postać $Bj=Lf=[s]-[t]$.
Zmiana znaku nie wpływa na energię prądu.

\begin{definicja}[Efektywna rezystancja między wierzchołkami]
Niech $G=(V,E)$ będzie skończonym spójnym grafem prostym,
$|V|\ge2$, z dodatnimi wagami $w_e$. Niech $s,t\in V$
będą różnymi wierzchołkami, a $b=[s]-[t]\in\mathbb R^V$.
Wybieramy dowolne rozwiązanie równania $Lf=b$,
gdzie $L=BWB^T$ jest ważonym laplasjanem.
Liczbę $R_{\mathrm{eff}}(s,t):=f(s)-f(t)$ nazywamy
\emph{efektywną rezystancją} między $s$ i $t$.
\end{definicja}

\begin{stwierdzenie}[Istnienie rozwiązania i dodatniość oporu]
Niech $G=(V,E)$ będzie skończonym spójnym grafem
prostym z dodatnimi przewodnościami $w_e$,
$L$ jego laplasjanem ważonym, a $s,t\in V$ różnymi
wierzchołkami. Dla $b=[s]-[t]$ równanie $Lf=b$ ma
rozwiązanie, a każde dwa rozwiązania różnią się funkcją stałą.
Liczba $R_{\mathrm{eff}}(s,t)$ nie zależy od wyboru rozwiązania
i jest ściśle dodatnia. Ponadto
\[
R_{\mathrm{eff}}(s,t)=f^TLf
   =\sum_{uv\in E}w_{uv}(f(u)-f(v))^2.
\]
\end{stwierdzenie}

\begin{proof}
Ponieważ graf jest spójny, $\ker L$ jest prostą
rozpiętą przez wektor stały $\mathbf1$. Jako że $L$
jest symetryczna, mamy
$\operatorname{im}L=(\ker L)^\perp
=\{x\in\mathbb R^V:\sum_vx_v=0\}$.
Wektor $b=[s]-[t]$ ma sumę współrzędnych równą zero,
więc należy do obrazu $L$; rozwiązanie $f$ istnieje.
Dwa rozwiązania różnią się elementem $\ker L$,
czyli stałą, której dodanie nie zmienia $f(s)-f(t)$.

Mnożąc równanie $Lf=b$ skalarnie przez $f$,
otrzymujemy $f^TLf=f^Tb=f(s)-f(t)$.
Z twierdzenia \ref{thm:tg-energy} jest to suma
nieujemnych składników wskazana w tezie.
Nie może wynosić zero: wtedy $f$ byłaby stała,
a więc $Lf=0$, wbrew $b\ne0$. W konsekwencji
$R_{\mathrm{eff}}(s,t)>0$.
\end{proof}

\begin{samepage}
\begin{wniosek}[Rezystancja jako iloraz wag lasów i drzew]
\label{cor:tg-resistance-forest}
Niech $G=(V,E)$ będzie skończonym spójnym grafem
prostym z dodatnimi przewodnościami $w_e$,
a $s,t\in V$ różnymi wierzchołkami.
Niech $\tau_w(G)=\sum_T\prod_{e\in E(T)}w_e$,
gdzie suma przebiega po drzewach rozpinających.
Wówczas
\[
R_{\mathrm{eff}}(s,t)=
\frac{\displaystyle\sum_{F\in\mathcal F_{\{s,t\}}}
                \prod_{e\in E(F)}w_e}{\tau_w(G)}.
\]
W liczniku występują lasy rozpinające o dwóch
składowych, rozdzielających $s$ i $t$.
\end{wniosek}
\end{samepage}
\begin{proof}
Wybierzmy rozwiązanie $Lf=[s]-[t]$ z normalizacją
$f(t)=0$. Niech $J=V\setminus\{t\}$.
Ograniczenie układu do wierszy $J$ daje
$L_{JJ}f_J=[s]_J$. Macierz $L_{JJ}$ jest odwracalna
na mocy twierdzenia \ref{thm:tg-dirichlet} (lub
twierdzenia Kirchhoffa).
Wzór na element diagonalny macierzy odwrotnej daje
\[
f(s)=(L_{JJ}^{-1})_{ss}
=\frac{\det L_{V\setminus\{s,t\},V\setminus\{s,t\}}}
       {\det L_{JJ}}.
\]
W mianowniku stosujemy twierdzenie \ref{thm:tg-kirchhoff},
w liczniku twierdzenie \ref{thm:tg-rooted-forest-minor}.
Ponieważ $f(t)=0$, lewa strona jest rezystancją
efektywną, co dowodzi wzoru, także dla $|V|=2$.
\end{proof}

\begin{twierdzenie}[Zasada Thomsona: minimalna energia przepływu]
Niech $G=(V,E)$ będzie skończonym spójnym grafem
prostym z dodatnimi przewodnościami $w_e$, $B$ jego
macierzą incydencji, $W=\operatorname{diag}(w_e)$,
a $s,t\in V$ różnymi wierzchołkami. Niech $b=[s]-[t]$
i $L=BWB^T$. Przez
$\mathcal J:=\{j\in\mathbb R^E:Bj=b\}$ oznaczmy
zbiór wszystkich prądów krawędziowych realizujących
jednostkowy przepływ między $s$ i $t$.
Wtedy
\[
R_{\mathrm{eff}}(s,t)
=\min_{j\in\mathcal J}\sum_{e\in E}\frac{j_e^2}{w_e}.
\]
Minimum jest osiągane dokładnie przez prąd
$j_0=WB^Tf$, gdzie $f$ jest dowolnym rozwiązaniem $Lf=b$.
\end{twierdzenie}

\begin{proof}
Dla $j_0=WB^Tf$ mamy
$Bj_0=BWB^Tf=Lf=b$, więc $j_0\in\mathcal J$.
Każdy inny prąd $j\in\mathcal J$ możemy jednoznacznie
zapisać jako $j=j_0+h$ z $Bh=0$.
Rozwijając formę kwadratową energii,
otrzymujemy
\[
\begin{aligned}
 j^TW^{-1}j
 &=j_0^TW^{-1}j_0+2h^TW^{-1}j_0+h^TW^{-1}h\\
 &=f^TBWB^Tf+2h^TB^Tf+h^TW^{-1}h\\
 &=R_{\mathrm{eff}}(s,t)+h^TW^{-1}h.
\end{aligned}
\]
W drugim przejściu użyliśmy $j_0=WB^Tf$,
a w trzecim $h^TB^Tf=(Bh)^Tf=0$ i tożsamości
$R_{\mathrm{eff}}(s,t)=f^TLf$.
Ponieważ $W^{-1}$ jest dodatnio określona,
ostatni składnik jest nieujemny i znika wyłącznie
dla $h=0$. Zatem $j_0$ jest jedynym minimizatorem,
a wartość minimum wynosi $R_{\mathrm{eff}}(s,t)$.
\end{proof}

\begin{twierdzenie}[Monotoniczność Rayleigha]
\label{thm:tg-resistance-rayleigh}
Niech $G=(V,E)$ będzie skończonym spójnym grafem
prostym, $s,t\in V$, $s\ne t$, a $w,w'$ dwiema
rodzinami dodatnich przewodności na $E$, takimi że
$w'_e\ge w_e$ dla każdej krawędzi $e$.
Niech $R_w$ i $R_{w'}$ oznaczają rezystancje
efektywne dla odpowiednich przewodności.
Wtedy $R_{w'}(s,t)\le R_w(s,t)$.
Równość zachodzi dokładnie wtedy, gdy prąd minimalizujący
energię dla $w$ znika na każdej krawędzi, której
przewodność ściśle zwiększono.
\end{twierdzenie}

\begin{proof}
Niech $j$ będzie jedynym minimizatorem dla $w$.
Zbiór dopuszczalnych prądów $Bj=[s]-[t]$
nie zależy od przewodności. Z zasady Thomsona mamy
\[
R_{w'}(s,t)\le \sum_e\frac{j_e^2}{w'_e}
\le \sum_e\frac{j_e^2}{w_e}=R_w(s,t).
\]
Jeżeli zachodzi równość skrajnych wyrażeń, to
$\sum_e j_e^2(1/w_e-1/w'_e)=0$.
Każdy składnik jest nieujemny, a współczynnik
jest dodatni dokładnie dla $w'_e>w_e$.
Stąd na tych krawędziach $j_e=0$.

Odwrotnie, załóżmy, że $j$ znika na wszystkich
krawędziach o zwiększonej przewodności.
Niech $f$ spełnia $BWB^Tf=[s]-[t]$.
Z zasady Thomsona $j=WB^Tf$. Na krawędziach
o zmienionej wadze mamy $(B^Tf)_e=0$, gdyż
$w_e>0$ i $j_e=0$. Na pozostałych wag nie zmieniono.
W konsekwencji $W'B^Tf=WB^Tf=j$, a więc
$BW'B^Tf=[s]-[t]$. Ta sama funkcja $f$ rozwiązuje
równanie dla nowych przewodności. Różnica
$f(s)-f(t)$ pozostaje bez zmian, co daje równość.
\end{proof}

\begin{wniosek}[Dodawanie krawędzi]
Niech $G$ będzie skończonym spójnym grafem prostym
z dodatnimi przewodnościami, a $s,t$ jego różnymi
wierzchołkami, a $L$ jego ważonym laplasjanem.
Dodanie krawędzi między różnymi niesąsiadującymi
wierzchołkami $u,v$ z dodatnią przewodnością nie zwiększa
$R_{\mathrm{eff}}(s,t)$. Zmniejsza ją ściśle wtedy i tylko
wtedy, gdy $f(u)\ne f(v)$ dla rozwiązania $Lf=[s]-[t]$
w pierwotnym grafie.
\end{wniosek}

\begin{proof}
Przedłużamy stary prąd minimalizujący przez zero na
nowej krawędzi. Jest on dopuszczalny i ma tę samą
energię, więc nowe minimum nie przekracza starego.
Jeśli minima są równe, przedłużony prąd jest
minimizatorem także w nowym grafie. Z zasady Thomsona
pochodzi z potencjału $g$, a zerowy prąd na nowej
krawędzi daje $g(u)=g(v)$. Zachowujemy pomocnicze orientacje starych krawędzi.
Prądy pochodzące z $f$ i $g$ są na nich identyczne,
czyli $WB^T(f-g)=0$. Odwracalność $W$ daje
$B^T(f-g)=0$. Spójność daje stałość $f-g$,
a stąd $f(u)=f(v)$.
Odwrotnie, jeśli $f(u)=f(v)$, dodatkowa krawędź
nie zmienia wartości laplasjanu zastosowanego do $f$.
Funkcja $f$ rozwiązuje zatem nowe równanie i
rezystancja pozostaje bez zmian.
\end{proof}

\begin{przyklad}[Krawędź między punktami o tym samym potencjale]
Rozważmy cykl o krawędziach $sa,at,tb,bs$ i
przewodnościach jednostkowych. Funkcja
$f(s)=1$, $f(t)=0$, $f(a)=f(b)=1/2$ spełnia
$Lf=[s]-[t]$: w $s$ suma różnic wynosi $1$,
w $t$ wynosi $-1$, a w $a,b$ znika.
Rezystancja między $s$ i $t$ wynosi $1$.
Dodanie krawędzi $ab$ o dowolnej dodatniej
przewodności nie zmienia tej rezystancji,
ponieważ jej końce mają ten sam potencjał.
Natomiast zwiększenie przewodności $sa$ zmniejsza
rezystancję ściśle, gdyż pierwotny prąd na tej
krawędzi ma moduł $1/2$.
\end{przyklad}

\begin{przyklad}[Dwa równoległe sposoby przejścia prądu]
Dla ścieżki $s-a-t$ o przewodnościach równych jeden
przyjmujemy $b=[s]-[t]$. Rozwiązaniem $Lf=b$
jest funkcja $f(s)=1$, $f(a)=0$, $f(t)=-1$.
Stąd $R_{\mathrm{eff}}(s,t)=1-(-1)=2$:
dwa opory jednostkowe połączone szeregowo sumują się.

Dla trójkąta $K_3$ z wierzchołkami $s,t,a$
i wagami jednostkowymi mamy $L=3I-J$.
Wektor $f=(1/3,-1/3,0)$ w kolejności $(s,t,a)$
spełnia $Lf=(1,-1,0)=b$.
Wynika stąd
$R_{\mathrm{eff}}(s,t)=2/3$.
Bezpośrednia krawędź $st$ stanowi drogę o oporze $1$,
a druga możliwość $s-a-t$ ma opór $2$;
połączenie równoległe daje $(1^{-1}+2^{-1})^{-1}=2/3$.
Dodanie krawędzi zmniejszyło opór efektywny.
\end{przyklad}

\begin{center}
\begin{tikzpicture}[scale=.97,line cap=round]
  \begin{scope}[shift={(-3.2,0)}]
    \node[tgBlueVertex] (s) at (-1.1,-.2) {$s$};
    \node[tgTealVertex] (a) at (0,1.05) {$a$};
    \node[tgRedVertex] (t) at (1.1,-.2) {$t$};
    \draw[tgTree,line width=2.6pt] (s)--(a)--(t);
    \node[font=\footnotesize,text=tgNavy] at (0,-.82) {$R_{st}=2$};
  \end{scope}
  \begin{scope}[shift={(3.2,0)}]
    \node[tgBlueVertex] (s) at (-1.1,-.2) {$s$};
    \node[tgTealVertex] (a) at (0,1.05) {$a$};
    \node[tgRedVertex] (t) at (1.1,-.2) {$t$};
    \draw[tgTree,line width=2.6pt] (s)--(a)--(t);
    \draw[tgAccent,line width=2.6pt] (s)--(t);
    \node[font=\footnotesize,text=tgNavy] at (0,-.82) {$R_{st}=2/3$};
  \end{scope}
\end{tikzpicture}

{\footnotesize Dodanie czerwonej krawędzi jednostkowej
zwiększa przewodność pomiędzy $s$ i $t$.}
\end{center}


\newpage
\section{Luka spektralna i geometryczne znaczenie laplasjanu}
\label{sec:tg-spectral-gap}

\begin{definicja}[Druga wartość własna]
Niech $G$ będzie spójnym grafem prostym o $n\ge2$
wierzchołkach i laplasjanie nieważonym $L$.
Liczbę $\lambda_2(G)$, czyli drugą w rosnącym uporządkowaniu
wartości własnych $L$ z uwzględnieniem krotności, nazywamy \emph{spójnością algebraiczną}
(lub \emph{luką spektralną} laplasjanu).
\end{definicja}

\begin{twierdzenie}[Zasada Rayleigha]
Niech $G=(V,E)$ będzie skończonym spójnym grafem
prostym o $n\ge2$ wierzchołkach, z jednostkowymi
wagami. Niech $\lambda_2(G)$ oznacza drugą najmniejszą
wartość własną jego laplasjanu, liczoną z krotnościami.
Wówczas
\[
\lambda_2(G)=\min_{\substack{x\in\mathbb R^V\setminus\{0\}\\
               \sum_{v\in V}x_v=0}}
\frac{\sum_{uv\in E}(x_u-x_v)^2}{\sum_{v\in V}x_v^2}>0.
\]
\end{twierdzenie}

\begin{proof}
Macierz $L$ jest rzeczywista i symetryczna, więc posiada
ortonormalną bazę wektorów własnych. Wektor
$\mathbf1/\sqrt n$ odpowiada wartości własnej zero.
Każdy $x\perp\mathbf1$ możemy zapisać w pozostałej
części tej bazy jako $x=\sum_{j=2}^na_ju_j$.
Wtedy $x^TLx=\sum_{j=2}^n\lambda_ja_j^2$,
a $\|x\|^2=\sum_{j=2}^na_j^2$.
Iloraz Rayleigha jest więc średnią ważoną liczb
$\lambda_2,\ldots,\lambda_n$ o nieujemnych wagach,
a minimum przyjmuje na wektorze własnym dla
$\lambda_2$. Związek z sumą po krawędziach wynika
z twierdzenia \ref{thm:tg-energy}.
W spójnym grafie jedynymi wektorami zerowej energii
są stałe. W przestrzeni prostopadłej do $\mathbf1$
nie ma niezerowego takiego wektora, stąd $\lambda_2>0$.
\end{proof}

Liczba $\lambda_2$ mierzy opór wobec zmian funkcji między
połączonymi wierzchołkami. W następnym etapie, po wprowadzeniu
grafów planarnych i zanurzeń grafów w powierzchniach,
połączymy laplasjan z topologią komórek wyższych wymiarów,
kompleksami dualnymi i dalszymi przykładami dyskretnej geometrii.

% ============================================================
% Iteracja 4, rozdział 9
% ============================================================
